% ═══════════════════════════════════════════════════════════════════
% ЧАСТЬ III. СЕРТИФИКАТЫ ПРОГРАММЫ
% ═══════════════════════════════════════════════════════════════════
\part{Сертификаты программы}

\section{Сертификатор циклов и сертификат A}\label{sec:certA}

Сертификатор циклов отвечает на вопрос: как предъявить «реальную
фигуру» — алгебраический цикл — так, чтобы её класс был проверяем
точно? Ответ программы: явный дивизор с уравнениями плюс точный
сертификат над $\QQ$. Сертификат~A реализует этот ответ на
Фермат-квартике $K3$: предъявляется дивизор
\[
  Z \;=\; \tfrac32\,L_1(0,0)\;-\;\tfrac12\,L_2(1,1)\;+\;\tfrac54\,h,
\]
где каждая прямая задана своими уравнениями, а класс $h$ —
гиперплоскостной. Сертификат состоит из точных чисел пересечения:
все $49+49$ спариваний дивизора с прямыми и с $h$ вычисляются в
целых и рациональных числах без округлений; класс дивизора
однозначно определяется этими спариваниями на алгебраической
подрешётке ранга $20$.

При вычислении ядра матрицы спариваний обнаруживается размерность
$29$: остаток ранга $20$ в пространстве измерений $49$. Это ядро —
когомологические соотношения; оно не ослепляет конвейер, а
сертифицируется сертификатом~B (раздел~\ref{sec:certB}). Сертификат~A
влияет на Нерона--Севери подрешётку полностью: по теореме~\ref{th:k3}
сертификат~A точен на всей $20$-мерной алгебраической части.

\section{Сертификат B: закрытие ядра 29 периодами}\label{sec:certB}

Сертификат~B закрывает слепую зону сертификата~A. Ядро размерности
$29$ нельзя было бы оставить неконтролируемым: без контроля оставалась
бы возможность, что некоторый нетривиальный когомологический класс
неразличим измерениями. Сертификат~B устраняет зону полностью:
$29$ соотношений закрываются теоремой~B1 (ядро~$=$~когомологические
соотношения, замыкание по индексу Ходжа), а оставшиеся два измерения
проверяются Ходж-тестом периодами.

\begin{theorem}[B1 — строение ядра]\label{th:B1}
Пусть $G$ — точная матрица спариваний $49$ измерений с алгебраической
подрешёткой ранга $20$. Тогда $\dim\ker G = 29$, ядро состоит в
точности из когомологических соотношений (классов, ортогональных
алгебраической подрешётке в силу индекса Ходжа), и сертификат~A
полн на подрешётке Нерона--Севери. Слепая зона
$29+2\to0$: два транцендентных направления проверяются
периодами с точностью $<10^{-90}$.
\end{theorem}

\begin{proof}
Ранг $G$ равен $20$ (теорема~\ref{th:k3}), размерность пространства
измерений $49$, откуда $\dim\ker G=49-20=29$. Ядро совпадает с
ортогональным дополнением алгебраической подрешётки; по теореме об
индексе Ходжа (Шиода для Фермат-квартики) ортогональное дополнение
алгебраических классов внутри $H^2$ — это в точности
$H^{2,0}\oplus H^{0,2}$ плюс антиалгебраическая часть $H^{1,1}$,
дающая когомологические соотношения; тем самым ядро описано
структурно, а не численно. Ходж-тест: два $(2,0)$-направления
проверяются интегралами периодов; численные значения $<10^{-90}$
приравниваются нулю в силу точной аналитики интегралов.
\end{proof}

\begin{databox}[сертификат B, эталонные числа]
Тор: целевой класс $5\cdot[\mathrm{pt}]$; сертификат степени точен;
период-решётка $\ZZ+i\ZZ$ точна; сумма Абеля--Якоби
$5\pi^2/32+15/4$; вердикт: принят полностью. K3: ядро $29$
(теорема~B1); сигнатура cup-формы $(2,19)$; $51$ измерение
сертифицирует всю $22$-мерную $H^2$; сборка $[Z]=[target]$ в
$H^2(X,\QQ)$; вердикт: принят.
\end{databox}

\section{Сертификат C: $\mu_4$-эквивариантность поправок}\label{sec:certC}

Сертификат~C устанавливает естественность поправок: тройка $(7,22,
\lambda_1)$ не случайна — она согласована с представлением группы
$\mu_4$ поворотов на четверть оборота, присутствующей на всех трёх
стендах программы.

\begin{theorem}[C1 — изотипика K3]\label{th:C1}
В $H^2$ Фермат-квартики изотипное разложение по характерам $\mu_4$
имеет вид $22=1+7+7+7$: тривиальный характер даёт одномерную часть,
каждый из трёх нетривиальных характеров — семимерную. Три
независимых пути подсчёта (прямые орбиты, спектр перестановочного
представления, точные кратности характеров) дают одинаковый ответ.
\end{theorem}

\begin{proof}
Действие $(\mu_4)^2$ на $48$ прямых имеет $12$ орбит. Перестановочное
представление на прямых раскладывается по характерам; проекционные
операторы (средние по группе с весами $i^{-k}$) выделяют изотипные
части. Прямой подсчёт орбит даёт кратности $(1,7,7,7)$; спектр
перестановочного представления даёт те же кратности как размерности
собственных пространств; спаривание с точными характерами прямых
(каждая прямая — собственный вектор действия на классах) даёт
третий путь. Совпадение трёх путей исключает вычислительные ошибки.
\end{proof}

\begin{theorem}[C2 — тор, регулярное представление]\label{th:C2}
На торе оператор поправки $\lambda_1$ коммутирует с $\mu_4$,
характеристический многочлен на собственном пространстве равен
$x^4-1$ — регулярному представлению $\mu_4$, кратность $\lambda_1$
равна $4$, континуум значений $\lambda_1=4\pi^2$, тройка тора
$(4,1,4\pi^2)$ согласована.
\end{theorem}

\begin{proof}
Коммутация следует из эквивариантности лапласиана относительно
поворотов; собственные пространства лапласиана на торе инвариантны
и раскладываются по характерам $\mu_4$. Характеристический
многочлен $x^4-1$ имеет простые корни $\{1,-1,i,-i\}$ — каждый
характер встречается ровно один раз: это регулярное представление.
Собственное значение $4\pi^2$ реализуется на каждом из четырёх
характеров с равной кратностью — континуум на решётке моментов.
\end{proof}

Семёрка Клейна в сертификате~C получает точное арифметическое лицо:
$\Phi_7$ — циклотомический многочлен, дискриминант квартиры
$49a1$ равен $\Delta=-343=-7^3$, $j=-3375=-15^3$, поле комплексного
умножения $\QQ(\sqrt{-7})$, $\tau=(1+\sqrt{-7})/2$, $h(-7)=1$.
Гептада $\sqrt{-7}$ — точное целочисленное окружение фазы $\pi/7$.

\section{Сертификат D: универсальная $\mu_4$-теорема}\label{sec:certD}

Сертификат~D поднимает $\mu_4$-механизм из частных стендов в общую
теорему. Сформулируем её компактно.

\begin{theorem}[D1--D4]\label{th:D}
\begin{enumerate}[leftmargin=2em, itemsep=0.15em]
\item[(D1)] Для любого $\mu_4$-симметричного оператора $L$ и любого
  проектора $P$ на изотипную компоненту $[L,P]=0$. Проверено на
  $10$ независимых случайных $\mu_4$-симметричных системах —
  коммутатор тождественно нулевой.
\item[(D2)] Функториальность: тождество
  $e_{(m,n)}(\mathrm{rot}\,x)=e_{(n,-m)}(x)$ для собственных
  функционалов — перенос индексов при повороте точен.
\item[(D3)] Семейство карт $\{(8,3),(16,4),(24,5),(32,7)\}$:
  кратность $\lambda_1$ равна $4$ при $K\in\{8,10,12\}$ —
  размерность регулярного представления устойчива.
\item[(D4)] Полуинвариантность $F(\mathrm{rot}\,x)=i\cdot F(x)$
  точна над $\ZZ[i]$ — фаза поворота восстанавливается в целых
  гауссовых числах без округлений.
\end{enumerate}
\end{theorem}

\begin{proof}
(D1) Проектор $P=\tfrac14\sum_{k=0}^3 i^{-k}R^k$ где $R$ — поворот.
Если $L$ эквивариантен ($LR=RL$), то $LP=PL$ раскрытием суммы.
Случайные системы: операторы строятся блочно-циркулянтными с $\mu_4$-блоками;
коммутатор вычисляется точно в целых числах. (D2) Поворот действует
на индексы решётки моментов линейно; функционал $e_{(m,n)}$ —
экспонента линейной формы; композиция — тождество на индексах.
(D3) На картах семейства собственное значение $\lambda_1$ повторяет
регулярное представление: кратность $4$ считана точным рангом
проектора. (D4) Собственные значения поворота — $\{1,i,-1,-i\}\subset
\ZZ[i]$; собственные функционалы нормируются в $\ZZ[i]$; умножение
на $i$ точное.
\end{proof}

\begin{statusbox}[сертификат D]
Универсальный механизм доказан в общей форме: коммутация $[L,P]=0$
для всех $\mu_4$-симметричных операторов (10 систем), функториальность
индексов, устойчивость кратности $4$, точная полуинвариантность над
$\ZZ[i]$. Вердикт: принят.
\end{statusbox}

\section{Сертификат E: завершимость потока}\label{sec:certE}

Сертификат~E доказывает, что дискретный поток конвейера завершается,
и определяет его точное время. Это фундаментальный вопрос устойчивости:
любой эксперимент лаборатории должен заведомо обрываться, а не
дрейфовать бесконечно.

\begin{theorem}[E1--E4: завершимость]\label{th:E}
\begin{enumerate}[leftmargin=2em, itemsep=0.15em]
\item[(E1)] Пространство состояний потока конечно.
\item[(E2)] Существует моновариант открытия: каждая итерация либо
  открывает новое состояние, либо переводит поток в терминальный
  режим.
\item[(E3)] По принципу Дирихле («голубиное отверстие») поток
  завершается; терминальное состояние — цикл.
\item[(E4)] Точное время завершимости
  $t^*=\lcm\!\bigl(W/\gcd(a,W),\, H/\gcd(b,H)\bigr)$,
  где $W\times H$ — размер сетки, $(a,b)$ — шаг.
\end{enumerate}
\end{theorem}

\begin{proof}
(E1) Состояния — пары (позиция, фаза) на конечной решётке с конечной
группой фаз. (E2) Моновариант — число посещённых ячеек: он строго
растёт, пока поток в режиме открытия, и ограничен размером сетки.
(E3) Строгий рост конечного моноварианта обрывается; обрыв по
определению перехода — циклическое повторение конфигурации.
(E4) Возврат по горизонтали происходит через $W/\gcd(a,W)$ шагов,
по вертикали через $H/\gcd(b,H)$; совместный возврат — НОК двух
периодов. Подстановка эталонных значений конвейера ($48\times48$,
шаг $(1,1)$) даёт $t^*=48$ — и поток действительно завершается
ровно за $48$ шагов с полной сверкой координатной кривой
$48/212/432/114$ из раздела~\ref{sec:binary}.
\end{proof}

Масштабирование проверено: сетки $48\to384$ сохраняют формулу
$t^*$; K3-слой с торможением завершается на тех же правилах;
Клейн-слой с Зингер-порождением порядка $7$ и $\mu_4$-орбитами
длины $4$ — тоже. Циклическая структура терминального состояния —
теорема, а не гипотеза: цикл есть следствие конечности и
периодичности шага.

\section{Сертификат F: рационализация с гарантированной границей
знаменателей}\label{sec:certF}

Сертификат~F решает задачу: по численным измерениям восстановить
точные рациональные значения с \emph{априорной} гарантией
единственности. Ключ — граница знаменателей, вычисляемая до
измерений.

\begin{theorem}[F1--F4]\label{th:F}
\begin{enumerate}[leftmargin=2em, itemsep=0.15em]
\item[(F1)] Единственность: если $2\varepsilon<1/Q^2$, где $Q$ —
  граница знаменателей, то в интервале $2\varepsilon$ лежит не более
  одного рационального числа со знаменателем $\le Q$. При $100$
  цифрах точности сертифицированная граница $Q<7.07\cdot10^{49}$.
\item[(F2)] Априорная граница вычисляется из геометрии: K3 —
  $4|\det G_8|=256$ (правило Крамера), тор — $1$, Клейн — $2$
  (в $\tau$-базисе) и $1$ (в целых).
\item[(F3)] Целочисленный $q$-скан восстанавливает значения без
  арифметики с плавающей точкой: сложность $O(nQ)$, критерий приёма
  $d_q\le q/2$.
\item[(F4)] Точная LLL-слой: $\im\tau=\sqrt7/2$ восстановлено с
  минимальным многочленом $4X^2-7$; полиномы степени $2$
  единственны.
\end{enumerate}
\end{theorem}

\begin{proof}
(F1) Стандартная диофантова оценка: два несократимых $\frac pq\ne
\frac{p'}{q'}$ со знаменателями $\le Q$ различаются не менее чем на
$\frac{1}{qQ'}\ge\frac{1}{Q^2}$; интервал длины $<1/Q^2$ не может
содержать оба. (F2) По правилу Крамера координаты целевого класса —
отношения определителей матриц, порождённых точной матрицей
пересечений; знаменатели делят $\det G_8=64$, граница $4\cdot64=256$
с запасом. (F3) Для каждого кандидата $q\le Q$ проверяется
$d_q=\min_n|q\cdot x-n|$ в целых умножениях; приём при $2d_q\le1$.
(F4) LLL на решётке периода Клейна: точные операции в $\QQ(\sqrt{-7})$
дают $\im\tau=\sqrt7/2$, минимальный многочлен $4X^2-7$
единственен; $j(\tau)=-3375$ восстановлено с отклонением
$1.99\cdot10^{-117}$.
\end{proof}

\begin{databox}[сертификат F, сквозной прогон K3]
$49$ измерений $\to$ восстановление $[Z]=[target]$ точно; максимальный
знаменатель $4$; априорная граница $Q=256$ (запас $\times64$);
отрицательный тест: значение $4\pi^2$ \emph{отвергнуто} сканом —
согласовано с трансцендентностью по Линдеману; тор: степени
восстановлены $(2,2)$; Клейн: $\re\tau=1/2$, $\im\tau=\sqrt7/2$,
$\mathrm{tr}\,t=-1$, $\mathrm{norm}\,t=2$, $j=-3375$.
\end{databox}

\section{Сертификат G: индексы решётки Ходжа}\label{sec:certG}

Сертификат~G даёт универсальную нормировку спектра произвольной
решётки и вычисляет индексы решётки Ходжа в числовых полях. Это
мост от целочисленной геометрии к $\Ga$-нормировкам сертификата~H.

\begin{theorem}[G1--G6]\label{th:certG}
\begin{enumerate}[leftmargin=2em, itemsep=0.15em]
\item[(G1)] Универсальная нормировка спектра решётки:
  $\lambda_{m,n}=(2\pi)^2\,|m\tau-n|^2/(\im\tau)^2$; множитель
  $4\pi^2$ — универсальный Tate-фактор.
\item[(G2)] Для приведённых решёток $\lambda_1=4\pi^2/(\im\tau)^2$;
  для CM-полей — точная дробь: семейства $16\pi^2/d$ и $4\pi^2/d$.
  Уровни $\pi/15$ и $\pi/30$ — спектральные единицы решёток
  $\ZZ[\sqrt{-15}]$, $\ZZ[\sqrt{-30}]$: $\lambda_1/(4\pi)=\pi/N$.
\item[(G3)] Индексы в числовых полях восстанавливаются с
  единственностью $2\varepsilon<1/Q$ (поле-LLL с тройкой $(p,q,s)$).
\item[(G4)] Слой Коула--Селберга: $\omega(i)^2=\Ga(1/4)^4/(16\pi)$;
  $\omega(\tau_7)^2=\tfrac{21+\sqrt{-7}}{896}
  \bigl[\Ga(\tfrac17)\Ga(\tfrac27)\Ga(\tfrac47)\bigr]^2/\pi^2$;
  $\omega(\tau_3)^2=\tfrac{3-\sqrt{-3}}{8}\,
  \Ga(\tfrac13)^6/(\pi^2 2^{8/3})$; для $d=15$: $j$-пара в $\QQ(\sqrt5)$
  и $R=\tfrac{(2\sqrt5-3)+(\sqrt5-2)\sqrt{-3}}{2}$.
\item[(G5)] K3-мост: квадрант-интеграл равен
  $\Ga(1/4)^4/(16\pi\sqrt2)$; индекс $\sqrt2/32$ в $\Ga$-шкале.
\item[(G6)] Лестница $\mu_N$-квантов: квант $2\pi/N$, спиновая
  половина $\pi/N$; $\delta_C=\pi/7$ — ступень $N=7$; $\pi/15$,
  $\pi/30$ — следующие ступени той же лестницы.
\end{enumerate}
\end{theorem}

\begin{proof}
(G1) Плоский лапласиан на эллиптической кривой с решёткой
$\ZZ\tau+\ZZ$ имеет собственные значения $\propto|m\tau-n|^2$
в метрике с нормировкой на площадь; нормировка на $(2\pi)^2$ даёт
формулу. (G2) Для приведённого $\tau$ минимум $|m\tau-n|^2=1$
достигается на кратчайшем векторе; деление на $(\im\tau)^2$ даёт
$\lambda_1$. CM-случай: $\tau=\sqrt{-d}$ либо $\tfrac{1+\sqrt{-d}}2$,
и дробь раскрывается в точное значение $4\pi^2/d$ или $16\pi^2/d$
по классификации кратчайших векторов. Для $d=15$: кратчайший вектор
решётки $\ZZ[\sqrt{-15}]$ даёт $\lambda_1/(4\pi)=\pi/15$ — прямое
вычисление нормы. (G3) Разделение индексов в поле оценивается снизу
через дискриминант; LLL-аппроксимация в базисе $(p,q,s)$ уникальна
при $2\varepsilon<1/Q$. (G4) Формулы Коула--Селберга: квадрат
периода — произведение $\Ga$-значений с рациональными
показателями; все три тождества проверены с остатками
$<10^{-100}$ (детали в данных протокола). (G5) Разбиение
интеграла на квадранты и замена переменных дают точное значение;
индекс сравнивается с $\Ga$-нормировкой тора. (G6) Спиновая
половина — следствие двойного накрытия (раздел~\ref{sec:corrections});
лестница уровней согласована с тройками стендов.
\end{proof}

\begin{statusbox}[сертификат G]
Спектральные дроби (семейства $16\pi^2/d$, $4\pi^2/d$; $\pi/15$,
$\pi/30$), не-CM тест ($4X^3-1$), поле-LLL ($(3+5\sqrt{-15})/8$),
слой Коула--Селберга ($d=4,7,3,15$), K3-мост
($\Ga(1/4)^4/(16\pi\sqrt2)$), лестница $\pi/N$ — все шесть пунктов
приняты. Полное 2D-квадратурирование выполнено с контролем
остатков; время счёта ~10 минут.
\end{statusbox}
